\section{Moral Ecosystem Approach in AI Development}
\label{sec:5.4}
Urie Bronfenbrenner's ecological systems theory holds that a person develops inside several nested environments at once, from the immediate family outward to culture and historical time, each one shaping the individual while being reshaped by it in turn \autocite{bronfenbrenner1979ecology}. I borrow that structure here: an AI system's alignment with humane values is not set by its training data alone but by the interlocking systems around that training — the people who build it, the institutions they work inside, the culture and regulatory climate those institutions answer to, and the historical moment all of this happens to fall in.

Bronfenbrenner pictured these systems as concentric rings, each one shaping the ones inside it. The sections that follow move outward through the same structure, ring by ring, from the microsystem to the chronosystem.
